\section*{Infraestructura en la nube}
\phantomsection
\addcontentsline{toc}{section}{Infraestructura en la nube}
\label{sec:infraestructura}

La infraestructura en la nube constituye uno de los pilares fundamentales sobre los que se despliega y mantiene el sistema \textbf{Muncher}. Con el crecimiento de los servicios digitales y la creciente necesidad de entornos flexibles y escalables, la adopción del cómputo en la nube se ha convertido en un componente esencial del desarrollo moderno de software. Comprender cómo funciona la nube y qué ventajas ofrece frente a la infraestructura tradicional permite diseñar sistemas más robustos, seguros y fáciles de mantener.

En esta sección se abordarán los conceptos clave de la infraestructura de nube, diferenciando entre los enfoques tradicionales de aprovisionamiento dinámico y las prácticas modernas basadas en \textit{Infrastructure as Code} (IaC). Esta comparación permitirá entender cómo la automatización y la definición declarativa de recursos facilitan el control, la reproducibilidad y la trazabilidad de los entornos, aspectos críticos en la gestión del ciclo de vida de aplicaciones distribuidas.

Finalmente, se presentará un análisis de los principales proveedores de nube, evaluando sus características, fortalezas y limitaciones en relación con las necesidades específicas del proyecto. A partir de este estudio, se seleccionará la plataforma más adecuada para el despliegue de \textbf{Muncher}, detallando la estructura de infraestructura propuesta y justificando las decisiones técnicas adoptadas conforme a los requisitos de escalabilidad, integración continua y facilidad de mantenimiento.

\subsection*{Conceptos básicos de la nube e infraestructura}
\addcontentsline{toc}{subsection}{Conceptos básicos de la nube e infraestructura}

De acuerdo con la definición propuesta por el Instituto Nacional de Estándares y Tecnología (NIST, Mell y Grance, 2011) \cite{Grance2011}, la computación en la nube es un modelo que permite el acceso ubicuo, cómodo y bajo demanda a un conjunto compartido de recursos de computación configurables —como redes, servidores, almacenamiento, aplicaciones y servicios— que pueden aprovisionarse y liberarse rápidamente con un mínimo esfuerzo de gestión o interacción con el proveedor.

Este modelo se fundamenta en cinco características esenciales:

\begin{itemize}
	\item \textbf{Autoservicio bajo demanda}: los usuarios pueden obtener capacidades de cómputo, como tiempo de servidor o espacio de almacenamiento, automáticamente y sin intervención humana directa por parte del proveedor.

	\item \textbf{Acceso amplio a la red}: los recursos están disponibles mediante mecanismos estandarizados y pueden ser accedidos desde diversas plataformas cliente, como ordenadores, portátiles, tabletas o teléfonos móviles.

	\item \textbf{Agrupamiento de recursos} (\textit{resource pooling}): los proveedores aplican un modelo multiusuario en el que los recursos físicos y virtuales se asignan dinámicamente según la demanda, ofreciendo independencia respecto de la ubicación física del hardware.

	\item \textbf{Elasticidad rápida}: las capacidades pueden ampliarse o reducirse de manera inmediata —e incluso automática— conforme varía la carga de trabajo, proporcionando al usuario la sensación de contar con un conjunto de recursos prácticamente ilimitado.

	\item \textbf{Servicio medido}: el sistema controla, optimiza y registra el uso de los recursos mediante mecanismos de medición que garantizan transparencia tanto para el proveedor como para el consumidor.
\end{itemize}

En el contexto de este proyecto, la \textbf{infraestructura en la nube} se define como el conjunto de servicios y recursos proporcionados por el proveedor elegido que, en conjunto, conforman y sostienen el sistema \textbf{Muncher}. Estos elementos incluyen las capacidades de cómputo, almacenamiento, bases de datos, servicios de integración y herramientas de administración necesarias para que el sistema funcione de manera fiable y escalable en un entorno remoto.

A diferencia de una definición genérica, aquí la infraestructura se expresa en términos concretos del proveedor seleccionado. Esto significa que cada componente de la arquitectura se corresponde directamente con un servicio específico disponible en su catálogo. Por ejemplo, en el caso de Microsoft Azure, podrían emplearse \textit{App Services} para alojar aplicaciones web, \textit{Azure Functions} para ejecutar código sin servidor, \textit{Azure Storage} para gestionar datos y ficheros, o \textit{Azure Virtual Networks} para interconectar los servicios de forma segura. Si el proveedor elegido fuera Amazon Web Services (AWS), equivalentes funcionales serían recursos como \textit{AWS Lambda}, \textit{Amazon S3}, \textit{Amazon RDS}, o \textit{VPC}.

Definir la infraestructura con este enfoque permite que la arquitectura del sistema esté claramente alineada con las capacidades del proveedor, facilita la replicación del entorno en diferentes fases del proyecto (desarrollo, pruebas, producción) y asegura que todas las decisiones técnicas estén sustentadas por las herramientas y servicios concretos que la plataforma ofrece.


\subsection*{Infraestructura como código frente a infraestructura dinámica}
\addcontentsline{toc}{subsection}{Infraestructura como código frente a infraestructura dinámica}

El concepto de \textit{Infrastructure as Code} (IaC) se basa en aplicar prácticas propias del desarrollo de software a la gestión de sistemas e infraestructura. Esta metodología propone definir el entorno tecnológico mediante código, permitiendo que los procesos de provisión, configuración y actualización de los sistemas sean automáticos, repetibles y verificables.

Según Morris (2016) \cite{Morris2016}, la clave de IaC es tratar la infraestructura como si fuera código fuente, lo que hace posible aprovechar herramientas de control de versiones, pruebas automatizadas y orquestación de despliegues en su mantenimiento. De esta forma, se integran prácticas consolidadas del desarrollo ágil —como la Integración Continua (CI) y la Entrega Continua (CD)— en la administración de entornos en la nube.

Antes de adoptar el modelo de \textit{Infrastructure as Code} (IaC), es importante comprender las limitaciones que presenta la llamada infraestructura dinámica, un enfoque común en las primeras etapas de adopción del cómputo en la nube. Este modelo se basa en la provisión y configuración manual o semiautomatizada de servidores y entornos, normalmente mediante interfaces gráficas o scripts ad hoc. Aunque inicialmente resulta flexible y rápido para desplegar nuevos recursos, a medida que el sistema crece aparecen problemas de mantenimiento, consistencia y escalabilidad que afectan de forma directa la fiabilidad del entorno.

Uno de los principales riesgos es la proliferación de servidores y configuraciones difíciles de gestionar, fenómeno que Morris (2016) \cite{Morris2016} denomina “\textit{server sprawl}”. Cuando el equipo técnico puede crear recursos con facilidad pero sin control centralizado, el número de servidores aumenta más rápido que la capacidad de administrarlos correctamente. Esto genera entornos heterogéneos, con versiones de software y configuraciones distintas, donde los errores se multiplican y las tareas de actualización o parcheo se vuelven complejas y arriesgadas.

A esta diversidad no controlada se le suma la llamada "deriva de configuración" (\textit{configuration drift}), que ocurre cuando servidores idénticos en su origen se modifican con el tiempo de forma independiente. Pequeños ajustes, correcciones manuales o actualizaciones parciales provocan diferencias sutiles que terminan impidiendo reproducir el entorno completo, generando lo que Morris describe como “servidores copo de nieve” (\textit{snowflake servers}): máquinas únicas e irrepetibles que nadie se atreve a modificar por miedo a romper su funcionamiento. Cuando esta fragilidad se extiende a todo el ecosistema, se obtiene una “infraestructura Jenga”, donde cualquier cambio puede desestabilizar el conjunto.

En paralelo, la falta de estandarización genera lo que el autor denomina “miedo a la automatización” (\textit{automation fear}). Los equipos evitan ejecutar herramientas de automatización de manera continua, debido a la falta de confianza en los resultados o al temor de introducir inconsistencias en sistemas ya frágiles. Como consecuencia, las configuraciones se aplican de forma irregular y la automatización se limita a tareas puntuales, perpetuando un ciclo de desconfianza y errores.

Frente a estos problemas, el modelo de \textit{Infrastructure as Code} ofrece una alternativa estructurada y sostenible. Mediante la definición declarativa y versionada de la infraestructura, es posible reconstruir cualquier entorno desde cero de forma repetible, auditable y libre de dependencias ocultas. Cada cambio queda registrado en el control de versiones, integrándose con los flujos de integración y despliegue continuo. Este enfoque elimina la dependencia de configuraciones manuales, reduce la deriva entre entornos y evita la formación de servidores especiales no replicables.

Por estas razones, en el proyecto \textbf{Muncher} se opta por un modelo basado en IaC. Ello permite mantener coherencia entre entornos, asegurar la reproducibilidad del sistema y facilitar la automatización completa del proceso de despliegue. Además, proporciona una base sólida para escalar la infraestructura, incorporar nuevas funcionalidades y responder con agilidad ante cambios o incidencias en la nube.


\subsection*{Selección de un proveedor de nube}
\addcontentsline{toc}{subsection}{Selección de un proveedor de nube}

La elección de un proveedor de nube adecuado es un factor determinante en la arquitectura del sistema \textbf{Muncher}. La plataforma seleccionada debe ofrecer un equilibrio entre flexibilidad, coste y facilidad de gestión, permitiendo adaptar la infraestructura a las necesidades específicas del proyecto sin imponer restricciones excesivas derivadas de servicios demasiado abstractos o cerrados.

El sistema requiere una infraestructura orientada a desarrolladores y equipos técnicos que deseen mantener control total sobre los recursos desplegados. Se descartan soluciones que, aunque cómodas para proyectos pequeños, limitan la personalización o dificultan el acceso a configuraciones internas, como ocurre en plataformas tipo \textit{Netlify} o \textit{AWS Amplify}. El objetivo es disponer de un entorno que permita definir, conectar y ajustar los componentes del sistema de acuerdo con los requisitos del producto.

Los requisitos principales que se han establecido para la selección del proveedor son los siguientes:

\begin{itemize}
	\item \textbf{Alta capacidad de configuración}, evitando proveedores que limiten la personalización en favor de entornos “todo hecho”.
	\item \textbf{Compatibilidad con contenedores Docker}, para desplegar el back‑end de forma reproducible y aislada.
	\item \textbf{Soporte para aplicaciones web estáticas}, destinado al front‑end, con distribución eficiente de contenidos.
	\item \textbf{Herramientas de monitorización y observabilidad} sencillas, que permitan lectura de logs, obtención de métricas y seguimiento de errores.
	\item \textbf{Coste reducido}, con gasto mensual de unos pocos dólares, aprovechando capas gratuitas o planes económicos.
	\item \textbf{Bases de datos SQL gestionadas}, incluyendo replicación automática, copias de seguridad y mantenimiento simplificado.
\end{itemize}

Con los requisitos previamente establecidos, se ha realizado una evaluación detallada de los principales proveedores de servicios en la nube y de varias alternativas ligeras orientadas al despliegue simplificado. El objetivo es determinar qué plataformas se ajustan mejor a las necesidades de \textbf{Muncher} en términos de flexibilidad, capacidad de configuración, coste y facilidad de mantenimiento.

La comparación abarca tanto los tres grandes proveedores del mercado —\textbf{Amazon Web Services (AWS)}, \textbf{Microsoft Azure} y \textbf{Google Cloud Platform (GCP)}— como opciones más económicas y accesibles —\textbf{Heroku}, \textbf{Render} y \textbf{Vercel}— que pueden resultar especialmente adecuadas para entornos de desarrollo o fases tempranas del proyecto.

En las siguientes páginas se presentan las tablas individuales para cada proveedor, en las que se detallan sus características principales con respecto a los criterios definidos: capacidad de configuración, compatibilidad con contenedores, soporte de aplicaciones web estáticas, herramientas de observabilidad, coste estimado y disponibilidad de bases de datos SQL gestionadas.

% AWS
\begin{center}
	\renewcommand{\arraystretch}{1.2}
	\setlength{\tabcolsep}{6pt}
	\textbf{Amazon Web Services (AWS)}\\[4pt]
	\begin{tabularx}{\textwidth}{|l|X|}
		\hline
		\textbf{Criterio}               & \textbf{Descripción}                                                                                                                       \\
		\hline
		Alta capacidad de configuración & \checkmark\ Permite una configuración detallada y control total sobre recursos e infraestructura, sin restricciones en la personalización. \\
		\hline
		Soporte Docker / contenedores   & \checkmark\ Compatible mediante servicios como Amazon ECS y EKS para gestionar y orquestar contenedores Docker.                            \\
		\hline
		Web estáticas                   & S3 combinado con CloudFront para distribución global de contenido estático.                                                                \\
		\hline
		Monitorización y logs           & CloudWatch ofrece supervisión avanzada, alertas, métricas y visualización de registros en tiempo real.                                     \\
		\hline
		Coste mensual estimado mínimo   & \$5--10 (tier gratuito y escalado según consumo).                                                                                          \\
		\hline
		Bases de datos SQL gestionadas  & Amazon RDS y Aurora, con soporte para replicación, backup automático y mantenimiento gestionado.                                           \\
		\hline
	\end{tabularx}
\end{center}

\vspace{8pt}

% Azure
\begin{center}
	\renewcommand{\arraystretch}{1.2}
	\setlength{\tabcolsep}{6pt}
	\textbf{Microsoft Azure}\\[4pt]
	\begin{tabularx}{\textwidth}{|l|X|}
		\hline
		\textbf{Criterio}               & \textbf{Descripción}                                                                                                          \\
		\hline
		Alta capacidad de configuración & \checkmark\ Ofrece gran flexibilidad para definir recursos, conectarlos y gestionar infraestructura con scripts declarativos. \\
		\hline
		Soporte Docker / contenedores   & \checkmark\ Compatible con Azure Container Instances (ACI) y Azure Kubernetes Service (AKS).                                  \\
		\hline
		Web estáticas                   & Servicio de Static Web Apps, integrado con GitHub Actions para despliegue automático.                                         \\
		\hline
		Monitorización y logs           & Azure Monitor centraliza métricas, trazas y logs con herramientas complementarias como Application Insights.                  \\
		\hline
		Coste mensual estimado mínimo   & \$5--10 (nivel gratuito amplio, facturación por consumo).                                                                     \\
		\hline
		Bases de datos SQL gestionadas  & Azure SQL Database con alta disponibilidad, backup automático y escalado elástico.                                            \\
		\hline
	\end{tabularx}
\end{center}

\vspace{8pt}

% GCP
\begin{center}
	\renewcommand{\arraystretch}{1.2}
	\setlength{\tabcolsep}{6pt}
	\textbf{Google Cloud Platform (GCP)}\\[4pt]
	\begin{tabularx}{\textwidth}{|l|X|}
		\hline
		\textbf{Criterio}               & \textbf{Descripción}                                                                                                      \\
		\hline
		Alta capacidad de configuración & \checkmark\ Configuración completa de recursos, redes y permisos a nivel granular mediante la consola o herramientas CLI. \\
		\hline
		Soporte Docker / contenedores   & \checkmark\ Soporte nativo mediante Google Kubernetes Engine (GKE) y Cloud Run.                                           \\
		\hline
		Web estáticas                   & Alojamiento en Cloud Storage y distribución mediante Cloud CDN.                                                           \\
		\hline
		Monitorización y logs           & Cloud Logging y Cloud Monitoring permiten trazabilidad detallada del sistema y alerta automática.                         \\
		\hline
		Coste mensual estimado mínimo   & \$5--10 (plan gratuito con recursos limitados y facturación por uso real).                                                \\
		\hline
		Bases de datos SQL gestionadas  & Cloud SQL con soporte para PostgreSQL, MySQL y SQL Server gestionados automáticamente.                                    \\
		\hline
	\end{tabularx}
\end{center}

\vspace{8pt}

% Heroku
\begin{center}
	\renewcommand{\arraystretch}{1.2}
	\setlength{\tabcolsep}{6pt}
	\textbf{Heroku}\\[4pt]
	\begin{tabularx}{\textwidth}{|l|X|}
		\hline
		\textbf{Criterio}               & \textbf{Descripción}                                                                                                          \\
		\hline
		Alta capacidad de configuración & \xmark\ Plataforma altamente simplificada, con escasa personalización y control limitado sobre la infraestructura subyacente. \\
		\hline
		Soporte Docker / contenedores   & Parcial. Puede usar contenedores personalizados con Heroku Container Registry, aunque con limitaciones.                       \\
		\hline
		Web estáticas                   & Sí. Admite despliegues de aplicaciones estáticas a través de buildpacks o servidores ligeros.                                 \\
		\hline
		Monitorización y logs           & Monitorización y visualización de logs integradas en el panel, adecuadas para proyectos pequeños.                             \\
		\hline
		Coste mensual estimado mínimo   & \$0--7 (planes gratuito y Hobby adecuados para desarrollo o pruebas).                                                         \\
		\hline
		Bases de datos SQL gestionadas  & PostgreSQL nativo con copias de seguridad automáticas y mantenimiento gestionado.                                             \\
		\hline
	\end{tabularx}
\end{center}

\vspace{8pt}

% Render
\begin{center}
	\renewcommand{\arraystretch}{1.2}
	\setlength{\tabcolsep}{6pt}
	\textbf{Render}\\[4pt]
	\begin{tabularx}{\textwidth}{|l|X|}
		\hline
		\textbf{Criterio}               & \textbf{Descripción}                                                                                                                     \\
		\hline
		Alta capacidad de configuración & \checkmark\ Configuración flexible mediante archivos YAML, permitiendo definir servicios personalizados y dependencias.                  \\
		\hline
		Soporte Docker / contenedores   & \checkmark\ Soporte nativo a imágenes Docker sin necesidad de configuración avanzada.                                                    \\
		\hline
		Web estáticas                   & \checkmark\ Hosting de sitios estáticos gratuito, con integración continua desde GitHub o GitLab.                                        \\
		\hline
		Monitorización y logs           & Acceso sencillo a logs desde el panel y API; métricas básicas de rendimiento.                                                            \\
		\hline
		Coste mensual estimado mínimo   & \$0--7 (planes gratuitos y Starter, apropiados para prototipos o entornos pequeños).                                                     \\
		\hline
		Bases de datos SQL gestionadas  & PostgreSQL administrado con copias de seguridad automáticas y restauración rápida, sujeto a limitaciones relevantes en la capa gratuita. \\
		\hline
	\end{tabularx}
\end{center}

\vspace{8pt}

% Vercel
\begin{center}
	\renewcommand{\arraystretch}{1.2}
	\setlength{\tabcolsep}{6pt}
	\textbf{Vercel}\\[4pt]
	\begin{tabularx}{\textwidth}{|l|X|}
		\hline
		\textbf{Criterio}               & \textbf{Descripción}                                                                                                                 \\
		\hline
		Alta capacidad de configuración & \xmark\ Plataforma centrada en el despliegue automatizado de aplicaciones front-end, con opciones limitadas de configuración manual. \\
		\hline
		Soporte Docker / contenedores   & \xmark\ No admite ejecución directa de contenedores Docker.                                                                          \\
		\hline
		Web estáticas                   & \checkmark\ Altamente optimizado para sitios y aplicaciones front-end en React, Next.js y frameworks similares.                      \\
		\hline
		Monitorización y logs           & Sistema integrado de métricas y logs, con disponibilidad en el panel de control.                                                     \\
		\hline
		Coste mensual estimado mínimo   & \$0--7 (plan gratuito generoso con despliegues ilimitados para uso personal).                                                        \\
		\hline
		Bases de datos SQL gestionadas  & No directamente. Requiere uso de servicios externos como Supabase, Neon o PlanetScale para almacenamiento persistente.               \\
		\hline
	\end{tabularx}
\end{center}


Tras el análisis individual de los distintos proveedores, puede concluirse que los tres grandes —\textbf{AWS}, \textbf{Azure} y \textbf{GCP}— cumplen de forma plena los requisitos técnicos definidos para el sistema. Ofrecen una infraestructura madura y altamente configurable, con soporte completo para contenedores Docker, servicios de bases de datos SQL gestionadas y avanzadas herramientas de observabilidad. Sin embargo, su estructura de precios por consumo requiere conocer bien los recursos que se eligen y el coste de cada uno, para no incurrir en gastos innecesarios.

Las plataformas ligeras —\textbf{Heroku}, \textbf{Render} y \textbf{Vercel}— representan alternativas más sencillas y económicas. Su facilidad de uso y disponibilidad de planes gratuitos las hacen atractivas para proyectos en fases tempranas. No obstante, tanto Heroku como Vercel presentan ciertas limitaciones de flexibilidad y control sobre los recursos. En cambio, \textbf{Render} destaca como una solución intermedia: mantiene un modelo operativo simple, pero ofrece soporte nativo para contenedores Docker, alojamiento de aplicaciones estáticas, monitorización integrada y bases de datos SQL gestionadas. Además, su capa gratuita resulta atractiva para entornos de desarrollo o prototipado, si bien impone restricciones importantes en el caso de las bases de datos gestionadas, como se detalla a continuación.

Conviene matizar, no obstante, el alcance de la capa gratuita de \textbf{Render} en lo relativo a las bases de datos gestionadas. A diferencia de los servicios web, cuyo plan gratuito no tiene fecha de caducidad, las bases de datos PostgreSQL gratuitas disponen de 1 GB de almacenamiento fijo, no admiten ningún tipo de copia de seguridad, se limitan a una única instancia activa por espacio de trabajo y, sobre todo, expiran treinta días después de su creación, tras lo cual se conceden catorce días de gracia para migrar a un plan de pago antes de su eliminación definitiva (\url{https://render.com/docs/free}). La persistencia de los datos no puede sostenerse, por tanto, de forma indefinida sin coste alguno.

Alcanzado ese límite es necesario contratar una instancia de pago, en la que el almacenamiento se factura de manera independiente al cómputo. En ese escenario el gasto de la base de datos gestionada pasa a dominar el coste mensual del proyecto y puede igualar o superar el de una infraestructura equivalente desplegada sobre un proveedor generalista como \textbf{AWS} o \textbf{Azure}, donde el mismo servicio se dimensiona con mayor granularidad. Esta circunstancia convierte la eventual migración a otro proveedor en una decisión motivada por el coste, y no únicamente por la madurez técnica del sistema.

En el caso del proyecto \textbf{Muncher}, donde se busca conservar el control técnico total, disponer de soporte para contenedores, contar con bases de datos SQL gestionadas y mantener un coste mensual mínimo, la comparación conduce a seleccionar \textbf{Amazon Web Services}. Ofrece un entorno sólido y flexible, con ejecución de contenedores sobre Amazon ECS, bases de datos gestionadas mediante Amazon RDS y distribución del front-end a través de S3 y CloudFront; su facturación por consumo permite dimensionar cada componente de forma independiente y ajustar el gasto al uso real del sistema.

Frente a las plataformas ligeras, esta decisión evita la migración forzosa que impondría el agotamiento de la capa gratuita descrita anteriormente. Aunque \textbf{Render} resulte más sencillo de operar en las primeras fases del desarrollo, la caducidad de sus bases de datos gratuitas obligaría a trasladar la infraestructura a otro proveedor en un plazo corto, con el coste y el riesgo que conlleva rehacer el despliegue una vez el sistema está en funcionamiento. Construir los entornos de desarrollo y de producción sobre la misma infraestructura desde el inicio elimina ese trabajo y permite que las decisiones tomadas durante el desarrollo sigan siendo válidas en producción.


\subsection*{Elementos que componen la infraestructura}
\addcontentsline{toc}{subsection}{Elementos que componen la infraestructura}

Una infraestructura en la nube no se entiende únicamente como un conjunto de recursos aislados, sino como una arquitectura coherente donde cada componente cumple una función específica dentro del sistema y se integra con los demás para formar un entorno operativo completo. El propósito de esta sección es describir los distintos elementos que conforman la infraestructura seleccionada para el proyecto \textbf{Muncher}, analizando su papel en la arquitectura general y justificando las decisiones técnicas que motivan su adopción.

\subsubsection*{Front-end}
\addcontentsline{toc}{subsubsection}{Front-end}

Para el componente de \textbf{front‑end}, al tratarse de una aplicación de página única (\textit{Single Page Application}, SPA), fue necesario decidir cómo publicar los archivos resultantes de su compilación. \textbf{AWS} ofrece varias vías para ello: el alojamiento de sitios estáticos sobre \textbf{Amazon S3}, la combinación de dicho almacenamiento con la red de distribución \textbf{Amazon CloudFront}, o servicios gestionados de mayor nivel como \textit{AWS Amplify Hosting}. Las tres alternativas resultan válidas para servir contenido estático, pero difieren en el grado de control sobre la configuración, en el coste y en la capacidad de distribución geográfica.

Amazon S3 permite alojar un sitio web estático a partir de los archivos depositados en un \textit{bucket}, entendiendo por estático aquel cuyas páginas incluyen contenido fijo y, opcionalmente, scripts ejecutados en el cliente.\footnote{Documentación oficial: \url{https://docs.aws.amazon.com/AmazonS3/latest/userguide/WebsiteHosting.html}} Sobre este almacenamiento se sitúa \textbf{Amazon CloudFront}, que distribuye el contenido a través de una red mundial de centros de datos denominados \textit{edge locations}: cada petición se encamina a la ubicación que ofrece la menor latencia y, si el archivo no está en ella, se recupera del origen definido —en este caso el \textit{bucket} de S3— y se almacena en caché para las siguientes.\footnote{Documentación oficial: \url{https://docs.aws.amazon.com/AmazonCloudFront/latest/DeveloperGuide/Introduction.html}} CloudFront asigna además un dominio propio a cada distribución, admite el uso de un nombre de dominio alternativo y permite controlar el tiempo de permanencia de los archivos en caché. La transferencia de datos entre el origen y CloudFront no se factura cuando el origen es un servicio de AWS como S3, de modo que el coste se limita a la salida de datos hacia el usuario y al número de peticiones atendidas.

\begin{figure}[htbp]
	\centering
	\includegraphics[width=0.85\textwidth]{front-end-infrastructure.pdf}
	\caption{Infraestructura del front-end.}
	\label{fig:front-end-infrastructure}
\end{figure}

Se optó por la combinación de \textbf{S3 y CloudFront} porque el front‑end de \textbf{Muncher} no requiere ejecución de código en el servidor: basta con almacenar y servir los archivos del cliente compilado (HTML, CSS, JavaScript y recursos estáticos), lo que elimina la necesidad de instancias o procesos persistentes y reduce el coste de la infraestructura. Se descartó \textit{AWS Amplify Hosting} pese a ser la opción que la propia documentación de S3 recomienda para este caso, ya que su carácter de servicio totalmente gestionado abstrae buena parte de la configuración subyacente y entra en conflicto con el criterio de alta capacidad de configuración establecido al inicio de esta sección. Frente a ello, S3 y CloudFront ofrecen el mismo resultado —entrega global mediante CDN de una aplicación cliente— conservando el control explícito sobre el almacenamiento, la política de caché y la distribución.

\subsubsection*{API}
\addcontentsline{toc}{subsubsection}{API}

Para el componente de \textbf{API} la decisión es de naturaleza distinta a la del front-end: no se trata de una aplicación estática, sino que requiere cómputo para atender cada petición, de modo que la elección recae sobre el modelo de ejecución. \textbf{AWS} ofrece dos familias de opciones para ello. La primera son los contenedores gestionados —\textbf{Amazon ECS} sobre \textit{Fargate}, \textbf{AWS App Runner} o los servicios de contenedores de \textbf{Amazon Lightsail}—, en los que la aplicación se empaqueta en una imagen y se mantiene un proceso en ejecución de forma permanente. La segunda son las funciones de \textbf{AWS Lambda}, donde el proveedor ejecuta el código únicamente durante el tiempo que tarda en atender cada invocación. A esta elección se añade una segunda, independiente de la anterior: cómo exponer la API a Internet, para lo que caben \textbf{Amazon API Gateway}, un balanceador de carga de aplicación o la propia URL de función situada detrás de la distribución de CloudFront descrita en el apartado anterior.

El precio es lo que separa con mayor claridad a ambas familias. Un contenedor en ejecución continua se cobra por horas con independencia de que reciba tráfico: la tarea más pequeña que admite \textit{Fargate} —0,25 vCPU y 0,5 GB de memoria— supone, a los precios publicados para la región de Norte de Virginia, un gasto próximo a los 9 USD mensuales si permanece activa todo el mes,\footnote{Precios oficiales: \url{https://aws.amazon.com/ecs/pricing/}} y el plan más reducido de los servicios de contenedores de \textbf{Lightsail} asciende a 7 USD mensuales fijos.\footnote{Precios oficiales: \url{https://aws.amazon.com/lightsail/pricing/}} A ello se suma que \textbf{App Runner} dejó de admitir nuevos clientes el 30 de abril de 2026, por lo que ha dejado de ser una alternativa disponible.\footnote{Aviso oficial: \url{https://aws.amazon.com/apprunner/}} \textbf{Lambda}, por el contrario, no factura mientras no se invoca, y su capa gratuita —un millón de peticiones y 400.000 GB-segundo de cómputo al mes— es permanente y no expira a los doce meses como ocurre con otras.\footnote{Precios oficiales: \url{https://aws.amazon.com/lambda/pricing/}} Con el tráfico previsto para este sistema, el coste de la API es por tanto nulo, frente a un gasto fijo que por sí solo agotaría el margen que fija el criterio de coste reducido establecido al inicio de esta sección.

Se optó en consecuencia por \textbf{AWS Lambda}, y la segunda decisión se resolvió reutilizando la distribución de \textbf{CloudFront} que ya sirve el front-end. A ella se añadió un comportamiento de caché para las rutas \texttt{/api/*} cuyo origen es la URL de función, de manera que la API se publica en el mismo dominio que la interfaz y bajo esa ruta. Se descartó \textbf{API Gateway} porque su cometido —enrutar las peticiones y publicar la API— lo cumple ya la distribución existente, y porque su capa gratuita de un millón de llamadas cubre únicamente los doce primeros meses desde la creación de la cuenta,\footnote{Precios oficiales: \url{https://aws.amazon.com/api-gateway/pricing/}} de forma que incorporarlo equivaldría a programar un gasto futuro sin obtener ninguna capacidad que el sistema no tenga. El balanceador de carga se descartó por la misma razón que los contenedores: se factura por horas de existencia y no por uso.

\begin{figure}[htbp]
	\centering
	\includegraphics[width=0.85\textwidth]{api-infrastructure.pdf}
	\caption{Infraestructura de la API.}
	\label{fig:api-infrastructure}
\end{figure}

El modelo elegido impone dos limitaciones que conviene hacer explícitas. La primera es el arranque en frío: cuando una invocación no encuentra un entorno de ejecución ya disponible, al tiempo de respuesta se añade la inicialización del intérprete y la importación de las dependencias de la aplicación. La segunda afecta a las peticiones con cuerpo, ya que Lambda no admite cargas sin firmar cuando el origen está protegido de esta manera: los métodos POST y PUT exigen que el cliente calcule el SHA-256 del cuerpo y lo envíe en la cabecera \texttt{x-amz-content-sha256}.

\subsubsection*{Monitorización y observabilidad}
\addcontentsline{toc}{subsubsection}{Monitorización y observabilidad}

Para el componente de \textbf{monitorización} la decisión no versa sobre dónde se ejecuta el sistema, como en los dos apartados anteriores, sino sobre dónde se recogen, se almacenan y se consultan los datos que describen su funcionamiento. Se distinguen tres clases de datos, y los tres son necesarios para satisfacer los requerimientos de esta categoría: los \textit{logs}, que dejan escritos los sucesos que la aplicación considera relevantes; las trazas, que siguen una misma petición a través de los componentes que intervienen en atenderla y miden el tiempo consumido en cada uno; y las métricas, que agregan el comportamiento del sistema en series temporales.

\textbf{AWS} cubre este cometido con \textbf{Amazon CloudWatch} —que reúne el almacenamiento y la consulta de \textit{logs}, las métricas, los paneles y las alarmas— y con \textbf{AWS X-Ray}, que recoge las trazas. Fuera del proveedor existe un conjunto de plataformas de observabilidad cuyas capas gratuitas son considerablemente más amplias: \textbf{Grafana Cloud} ofrece 50 GB mensuales de \textit{logs} y otros tantos de trazas con catorce días de retención,\footnote{Precios oficiales: \url{https://grafana.com/pricing/}} \textbf{New Relic} concede 100 GB mensuales de datos de cualquier origen y un usuario con acceso completo,\footnote{Precios oficiales: \url{https://newrelic.com/pricing/free-tier}} y \textbf{Axiom} llega a 500 GB mensuales con treinta días de retención.\footnote{Precios oficiales: \url{https://axiom.co/pricing}} Una tercera vía consistiría en desplegar una solución propia, como \textbf{Grafana} junto a \textbf{Loki} y \textbf{Tempo}, sobre infraestructura del proyecto.

Esta última se descartó por la misma razón que condujo a elegir Lambda para la API: exige un proceso en ejecución permanente. El servidor más reducido capaz de sostener un sistema de este tipo consumiría por sí solo la totalidad del margen de gasto fijado para el proyecto, de modo que observar el sistema costaría más que ejecutarlo. Sostener una infraestructura que no factura mientras no se usa mediante otra que factura por horas anularía precisamente la ventaja que motivó la primera decisión.

Las plataformas externas se descartaron por dos motivos que se refuerzan entre sí. El primero es técnico: una función de Lambda no mantiene un proceso vivo entre invocaciones, por lo que enviar la telemetría a un destino externo obliga a hacerlo dentro de la propia petición —lo que añade latencia a cada respuesta— o a desplegar junto a la función un componente adicional que la recoja y la reenvíe. X-Ray no plantea ese problema, porque el entorno de ejecución de Lambda incluye un \textit{daemon} que recibe la telemetría de forma local y se encarga de enviarla.\footnote{Variables de entorno de Lambda: \url{https://docs.aws.amazon.com/lambda/latest/dg/configuration-envvars.html}} El segundo es económico: Lambda escribe sus \textit{logs} en CloudWatch sin que el destino sea configurable, de manera que ese volumen se consume en cualquier caso; remitirlo además a una plataforma externa supondría mantener dos sistemas de observabilidad para un único conjunto de datos.

Se optó en consecuencia por \textbf{CloudWatch} y \textbf{X-Ray}. Su capa gratuita cubre con holgura el tráfico previsto: 5 GB mensuales de datos de \textit{logs} —sumando la ingesta, el almacenamiento y el volumen examinado por las consultas—, 100.000 trazas registradas al mes, diez alarmas y tres paneles de hasta cincuenta métricas cada uno; superados esos límites, la ingesta se factura a 0,50 USD por GB y el almacenamiento a 0,03 USD por GB.\footnote{Precios oficiales: \url{https://aws.amazon.com/cloudwatch/pricing/}} Igual que se argumentó al descartar API Gateway, lo determinante no es la amplitud de la capa gratuita sino su permanencia: la de CloudWatch y X-Ray no expira a los doce meses de creada la cuenta, por lo que adoptarlas no programa un gasto futuro.

Elegido el destino, quedaba decidir cómo genera la API las trazas que lo alimentan. El rastreo activo de Lambda registra por sí solo cada invocación y separa la inicialización del entorno del tiempo de ejecución, pero no reparte este último entre las operaciones de la aplicación, que es lo que pide \reqref{RNF-MO-03}. La vía tradicional para obtener ese desglose, el SDK de X-Ray, entró en modo de mantenimiento en febrero de 2026 y dejará de recibir soporte en febrero de 2027, y \textbf{AWS} recomienda sustituirlo por \textbf{OpenTelemetry}, cuyas trazas X-Ray recibe sin cambios en la consola.\footnote{Calendario de fin de soporte del SDK de X-Ray: \url{https://docs.aws.amazon.com/xray/latest/devguide/xray-daemon-eos.html}} Adoptar en un proyecto nuevo una dependencia con fecha de retirada habría obligado a migrarla pocos meses después.

Se optó por el \textit{layer} de OpenTelemetry que publica \textbf{AWS} para Lambda, que aporta el SDK y lo configura para enviar las trazas al \textit{daemon} de X-Ray del entorno de ejecución.\footnote{Instrumentación de Lambda con OpenTelemetry: \url{https://docs.aws.amazon.com/xray/latest/devguide/migrate-xray-to-opentelemetry-python.html}} La aplicación no contiene instrumentación propia: \textbf{FastAPI} emite de forma nativa, desde su versión 0.142, un \textit{span} por petición con el nombre de la ruta que la atiende y otro por cada fase de su procesamiento —la resolución de dependencias, la ejecución de la operación y la serialización de la respuesta—, y el \textit{layer} los incorpora a la traza de la invocación. El precio de esta solución es el arranque en frío, que se alarga porque el SDK se importa antes que la aplicación.

\begin{figure}[htbp]
	\centering
	\includegraphics[width=0.85\textwidth]{monitoring-infrastructure.pdf}
	\caption{Infraestructura de la monitorización.}
	\label{fig:monitoring-infrastructure}
\end{figure}

El diseño satisface \reqref{RNF-MO-01} mediante el registro de los errores, \reqref{RNF-MO-03} mediante las trazas, que desglosan el tiempo de cada petición por operación, y de forma parcial \reqref{RNF-MO-04}, ya que el recuento de invocaciones mide peticiones atendidas y no usuarios activos. Queda sin satisfacer \reqref{RNF-MO-02}, por un motivo económico: la capa gratuita que cubriría las notificaciones es acotada y nada limita por sí solo el número de avisos que un sistema de alertas emite, de modo que permanecer dentro de ella exigiría añadir mecanismos de contención que son a su vez infraestructura capaz de fallar. Al tratarse de una prioridad \textit{Should Have}, la decisión es reversible y se revisará cuando el sistema reciba tráfico real.

El modelo presenta, por último, dos puntos ciegos que conviene hacer explícitos. El primero es CloudFront: la distribución no participa en las trazas, de modo que el tiempo transcurrido entre la petición del navegador y la invocación de la función no queda medido y la traza comienza en la función misma. El segundo es el propio navegador, ya que el comportamiento del front-end en el cliente —tiempos de carga y errores de JavaScript— no se recoge: el servicio de AWS destinado a ello, \textbf{CloudWatch RUM}, se factura a 1 USD por cada 100.000 eventos y su gratuidad se limita a una prueba inicial de un millón de eventos por cuenta, no a una asignación mensual recurrente.

